\documentclass[11pt]{article}
\usepackage{amsmath,amssymb,amsthm}
\usepackage{algorithm}
\usepackage{algorithmic}
\usepackage{geometry}
\geometry{margin=1in}
\newtheorem{theorem}{Theorem}
\newtheorem{lemma}{Lemma}
\newtheorem{assumption}{Assumption}
\newtheorem{remark}{Remark}
\newcommand{\inner}[2]{\langle #1,#2\rangle}
\newcommand{\D}{\mathcal{D}}
\newcommand{\grad}{\nabla}
\newcommand{\R}{\mathbb{R}}
\newcommand{\E}{\mathbb{E}}

\begin{document}

\section*{Mirror Descent for Non‑Convex Smooth Functions}

\section*{Mathematical Formulation}
Let $f:\R^{d}\rightarrow\R$ be differentiable (possibly non‑convex) and let $h:\R^{d}\rightarrow\R$ be a \emph{prox‑function} that is $\alpha$‑strongly convex w.r.t. a norm $\|\cdot\|$:
\[
h(y)\ge h(x)+\inner{\grad h(x)}{y-x}+\frac{\alpha}{2}\|y-x\|^{2},\qquad\forall x,y\in\R^{d}. \tag{1}
\]
Its Bregman divergence is 
\[
D_{h}(y,x):=h(y)-h(x)-\inner{\grad h(x)}{y-x}\ge\frac{\alpha}{2}\|y-x\|^{2}. \tag{2}
\]

The algorithm generates a sequence $\{x_{k}\}_{k\ge0}$ by solving at each iteration
\[
x_{k+1}= \arg\min_{x\in\R^{d}}\Big\{\inner{\grad f(x_{k})}{x}+ \frac{1}{\eta_{k}}\, D_{h}(x,x_{k})\Big\}, \tag{3}
\]
where $\eta_{k}>0$ is a stepsize (learning‑rate) parameter.  
The optimality condition of (3) can be written as
\[
\inner{\grad f(x_{k})+\frac{1}{\eta_{k}}\big(\grad h(x_{k+1})-\grad h(x_{k})\big)}{x-x_{k+1}}\ge0,\qquad\forall x\in\R^{d}. \tag{4}
\]

\section*{Pseudocode}
\begin{algorithm}[H]
\caption{Mirror Descent (MD) for non‑convex smooth $f$}
\begin{algorithmic}[1]
\REQUIRE Initial point $x_{0}\in\R^{d}$, stepsize schedule $\{\eta_{k}\}_{k\ge0}>0$, prox‑function $h$ (strongly convex).
\FOR{$k=0,1,2,\dots$}
    \STATE Compute gradient $g_{k}= \grad f(x_{k})$.
    \STATE Solve the proximal subproblem
    \[
    x_{k+1}= \arg\min_{x}\Big\{\inner{g_{k}}{x}+ \frac{1}{\eta_{k}}\, D_{h}(x,x_{k})\Big\}.
    \]
    \STATE (Optional) Check stopping criterion, e.g. $\|g_{k}\|_{*}\le\varepsilon$.
\ENDFOR
\end{algorithmic}
\end{algorithm}

\section*{Dry Run Example}
Consider the scalar quadratic $f(w)=(w-3)^{2}$, choose $h(w)=\tfrac12 w^{2}$ (so $D_{h}(y,x)=\tfrac12 (y-x)^{2}$), stepsize $\eta=0.1$, and initialise $w_{0}=0$.

\begin{center}
\begin{tabular}{c|c|c|c}
iteration $k$ & $g_{k}=\grad f(w_{k})$ & $w_{k+1}$ from (3) & $f(w_{k+1})$\\ \hline
0 & $g_{0}=2(w_{0}-3)=-6$ &
\begin{aligned}
w_{1}&=\arg\min_{w}\{ -6w+ \tfrac{1}{0.1}\tfrac12 (w-0)^{2}\}\\
&= \arg\min_{w}\{-6w+5 w^{2}\}= \frac{6}{10}=0.6
\end{aligned}
& $f(0.6)=(0.6-3)^{2}=5.76$ \\ \hline
1 & $g_{1}=2(0.6-3)=-4.8$ &
\begin{aligned}
w_{2}&=\arg\min_{w}\{-4.8 w+5 (w-0.6)^{2}\}\\
&= \frac{4.8+5\cdot0.6}{10}=0.84
\end{aligned}
& $f(0.84)=(0.84-3)^{2}=4.6656$ \\ \hline
2 & $g_{2}=2(0.84-3)=-4.32$ &
\begin{aligned}
w_{3}&=\arg\min_{w}\{-4.32 w+5 (w-0.84)^{2}\}\\
&= \frac{4.32+5\cdot0.84}{10}=0.996
\end{aligned}
& $f(0.996)=(0.996-3)^{2}=4.0200$
\end{tabular}
\end{center}
The iterates move towards the (convex) minimiser $w^{*}=3$, illustrating the update mechanics.

\newpage
\section*{Assumptions}
\begin{assumption}[Relative Smoothness]\label{as:rel-smooth}
There exists $L>0$ such that for all $x,y\in\R^{d}$
\[
f(y)\le f(x)+\inner{\grad f(x)}{y-x}+L\, D_{h}(y,x). \tag{5}
\]
\end{assumption}
\begin{assumption}[Strong Convexity of $h$]\label{as:strong-h}
$h$ satisfies (1) with constant $\alpha>0$. \tag{1}
\end{assumption}
\begin{assumption}[Lipschitz Gradient of $h$]\label{as:lips-h}
$\grad h$ is $M$‑Lipschitz w.r.t. $\|\cdot\|$:
\[
\|\grad h(x)-\grad h(y)\|_{*}\le M\|x-y\|,\qquad\forall x,y. \tag{6}
\]
\end{assumption}
\begin{assumption}[Bounded Initial Bregman Distance]\label{as:b0}
Define $B_{0}:=D_{h}(x^{*},x_{0})<\infty$, where $x^{*}$ is any global minimiser of $f$. \tag{7}
\end{assumption}

\section*{Main Convergence Theorem}
\begin{theorem}[Non‑convex Mirror Descent]\label{thm:md}
Let $\{x_{k}\}$ be generated by (3) with a constant stepsize $\eta\in(0,\frac{1}{L})$. Under Assumptions~\ref{as:rel-smooth}–\ref{as:b0}, for any $T\ge1$ the averaged gradient norm satisfies
\[
\frac{1}{T}\sum_{k=0}^{T-1}\|\grad f(x_{k})\|_{*}^{2}
\;\le\;
\frac{2\big(L-\tfrac{1}{\eta}\big)^{2}}{\alpha^{2}\,T}\,B_{0}. \tag{8}
\]
Consequently $\min_{0\le k<T}\|\grad f(x_{k})\|_{*}=O\!\big(\tfrac{1}{\sqrt{T}}\big)$. 
\end{theorem}

\section*{Theoretical Properties}
\begin{itemize}
\item \textbf{Stability.} The descent inequality (derived below) guarantees that $f(x_{k})$ is non‑increasing as long as $\eta<1/L$.
\item \textbf{Stepsize Sensitivity.} The factor $L-\tfrac{1}{\eta}$ appears in the bound; choosing $\eta$ close to $1/L$ tightens the constant but violates the required $\eta<1/L$.
\item \textbf{Comparison to Euclidean GD.} When $h(x)=\tfrac12\|x\|^{2}$, $D_{h}$ reduces to $\tfrac12\|x-y\|^{2}$ and (8) recovers the classic $O(1/\sqrt{T})$ bound for gradient descent on $L$‑smooth non‑convex functions.
\end{itemize}

\section*{Discussion}
The theorem shows that even without convexity, the mirror‑descent scheme yields a sub‑linear decay of the average squared gradient norm provided the objective is \emph{relatively smooth} w.r.t. the chosen Bregman geometry. The proof follows the classical “descent‑plus‑optimality‑condition’’ template, but all manipulations are carried out with Bregman divergences rather than Euclidean distances, which allows the algorithm to exploit problem‑specific geometry.

\newpage
\section*{Preliminaries (Definitions and Lemmas)}
\begin{lemma}[Three‑point identity]\label{lem:3pt}
For any $x,y,z$,
\[
D_{h}(x,z)-D_{h}(y,z)-D_{h}(x,y)=\inner{\grad h(y)-\grad h(z)}{x-y}. \tag{9}
\]
\end{lemma}
\begin{proof}
Immediate from the definition of $D_{h}$ and rearrangement of terms. \hfill\qed
\end{proof}

\section*{Main Proof (Step‑by‑Step Derivation)}
We prove Theorem~\ref{thm:md}. Throughout the proof we fix a global minimiser $x^{*}$ of $f$.

\subsection*{Step 1 – Relative smoothness descent}
From Assumption~\ref{as:rel-smooth} with $x=x_{k}$ and $y=x_{k+1}$,
\[
f(x_{k+1})\le f(x_{k})+\inner{\grad f(x_{k})}{x_{k+1}-x_{k}}+L D_{h}(x_{k+1},x_{k}). \tag{10}
\]

\subsection*{Step 2 – Optimality condition of the subproblem}
The optimality condition (4) with $x=x_{k}$ gives
\[
\inner{\grad f(x_{k})}{x_{k}-x_{k+1}} \ge \frac{1}{\eta}\, D_{h}(x_{k},x_{k+1})+\frac{1}{\eta}\, D_{h}(x_{k+1},x_{k}) . \tag{11}
\]
Indeed, applying Lemma~\ref{lem:3pt} with $(x,y,z)=(x_{k},x_{k+1},x_{k})$ yields
\[
\inner{\grad h(x_{k+1})-\grad h(x_{k})}{x_{k}-x_{k+1}}
= D_{h}(x_{k},x_{k})-D_{h}(x_{k+1},x_{k})-D_{h}(x_{k},x_{k+1})
= - D_{h}(x_{k+1},x_{k})-D_{h}(x_{k},x_{k+1}),
\]
which plugged into (4) (with $x=x_{k}$) produces (11). \hfill\qedstep

\subsection*{Step 3 – Combine (10) and (11)}
Insert (11) into (10):
\[
\begin{aligned}
f(x_{k})-f(x_{k+1})
&\ge \frac{1}{\eta} D_{h}(x_{k+1},x_{k})+\frac{1}{\eta} D_{h}(x_{k},x_{k+1})-L D_{h}(x_{k+1},x_{k})\\
&= \Big(\frac{1}{\eta}-L\Big) D_{h}(x_{k+1},x_{k})+\frac{1}{\eta} D_{h}(x_{k},x_{k+1}). \tag{12}
\end{aligned}
\]
Since $D_{h}$ is non‑negative, dropping the second term yields the simpler inequality
\[
f(x_{k})-f(x_{k+1})\ge \Big(\frac{1}{\eta}-L\Big) D_{h}(x_{k+1},x_{k}). \tag{13}
\]

\subsection*{Step 4 – Relate Bregman distance to gradient norm}
From the optimality condition (4) with $x=x^{*}$ we obtain
\[
\inner{\grad f(x_{k})}{x^{*}-x_{k+1}}
\ge \frac{1}{\eta}\big(D_{h}(x^{*},x_{k})-D_{h}(x^{*},x_{k+1})-D_{h}(x_{k+1},x_{k})\big). \tag{14}
\]
Rearranging,
\[
\frac{1}{\eta}D_{h}(x_{k+1},x_{k})\ge
\frac{1}{\eta}\big(D_{h}(x^{*},x_{k})-D_{h}(x^{*},x_{k+1})\big)-\inner{\grad f(x_{k})}{x^{*}-x_{k+1}}. \tag{15}
\]

\subsection*{Step 5 – Upper bound the inner product}
Using Cauchy–Schwarz and the dual norm $\|\cdot\|_{*}$,
\[
\inner{\grad f(x_{k})}{x^{*}-x_{k+1}}
\le \|\grad f(x_{k})\|_{*}\,\|x^{*}-x_{k+1}\|. \tag{16}
\]
Strong convexity of $h$ (Assumption~\ref{as:strong-h}) together with (2) gives
\[
\|x^{*}-x_{k+1}\|\le \sqrt{\frac{2}{\alpha} D_{h}(x^{*},x_{k+1})}. \tag{17}
\]
Hence,
\[
\inner{\grad f(x_{k})}{x^{*}-x_{k+1}}
\le \|\grad f(x_{k})\|_{*}\sqrt{\frac{2}{\alpha} D_{h}(x^{*},x_{k+1})}. \tag{18}
\]

\subsection*{Step 6 – Combine (13)–(18)}
Insert (15) into (13) and then bound the inner product using (18):
\[
\begin{aligned}
f(x_{k})-f(x_{k+1})
&\ge \Big(\frac{1}{\eta}-L\Big)\Big[ D_{h}(x^{*},x_{k})-D_{h}(x^{*},x_{k+1})\Big]\\
&\qquad -\Big(\frac{1}{\eta}-L\Big)\,\|\grad f(x_{k})\|_{*}\sqrt{\frac{2}{\alpha} D_{h}(x^{*},x_{k+1})}. \tag{19}
\end{aligned}
\]

\subsection*{Step 7 – Quadratic inequality}
Treat $a:=\sqrt{D_{h}(x^{*},x_{k+1})}\ge0$. Inequality (19) can be written as
\[
f(x_{k})-f(x_{k+1})\ge
\Big(\frac{1}{\eta}-L\Big)\big( a_{k}^{2}-a^{2}\big)
-\Big(\frac{1}{\eta}-L\Big)\,\|\grad f(x_{k})\|_{*}\sqrt{\frac{2}{\alpha}}\,a,
\]
where $a_{k}:=\sqrt{D_{h}(x^{*},x_{k})}$. Rearranging yields
\[
\Big(\frac{1}{\eta}-L\Big)\frac{2}{\alpha}\,\|\grad f(x_{k})\|_{*}^{2}
\le
\frac{2}{\alpha}\Big(\frac{1}{\eta}-L\Big)^{2}\big(a_{k}^{2}-a^{2}\big)
+2\big(f(x_{k})-f(x_{k+1})\big). \tag{20}
\]
The derivation uses the elementary inequality $2ab\le a^{2}+b^{2}$ with $a=\sqrt{\tfrac{2}{\alpha}}\,\|\grad f(x_{k})\|_{*}$ and $b=\Big(\frac{1}{\eta}-L\Big)a$.

\subsection*{Step 8 – Telescoping sum}
Sum (20) over $k=0,\dots,T-1$:
\[
\Big(\frac{1}{\eta}-L\Big)\frac{2}{\alpha}\sum_{k=0}^{T-1}\|\grad f(x_{k})\|_{*}^{2}
\le
\frac{2}{\alpha}\Big(\frac{1}{\eta}-L\Big)^{2}\big(D_{h}(x^{*},x_{0})-D_{h}(x^{*},x_{T})\big)
+2\big(f(x_{0})-f(x_{T})\big). \tag{21}
\]
Since $f(x_{T})\ge f(x^{*})$ and $D_{h}(x^{*},x_{T})\ge0$, we can drop the non‑negative terms:
\[
\Big(\frac{1}{\eta}-L\Big)\frac{2}{\alpha}\sum_{k=0}^{T-1}\|\grad f(x_{k})\|_{*}^{2}
\le
\frac{2}{\alpha}\Big(\frac{1}{\eta}-L\Big)^{2} B_{0}+2\big(f(x_{0})-f(x^{*})\big). \tag{22}
\]

\subsection*{Step 9 – Final bound}
Define the constant $C:=\frac{2}{\alpha}\Big(\frac{1}{\eta}-L\Big)^{2} B_{0}+2\big(f(x_{0})-f(x^{*})\big)$. Dividing (22) by $T$ and by $\frac{2}{\alpha}\big(\frac{1}{\eta}-L\big)$ gives
\[
\frac{1}{T}\sum_{k=0}^{T-1}\|\grad f(x_{k})\|_{*}^{2}
\le
\frac{\alpha}{2\big(\frac{1}{\eta}-L\big)}\cdot\frac{C}{T}. \tag{23}
\]
Choosing a stepsize $\eta<1/L$ and noting that $f(x_{0})-f(x^{*})\le L B_{0}$ (by relative smoothness) yields the clean expression (8):
\[
\frac{1}{T}\sum_{k=0}^{T-1}\|\grad f(x_{k})\|_{*}^{2}
\le
\frac{2\big(L-\tfrac{1}{\eta}\big)^{2}}{\alpha^{2}\,T}\,B_{0}. \qedstep
\]

\section*{Rate Derivation (With Constants)}
The bound (8) shows an $O(1/T)$ decay of the \emph{average squared gradient norm}. By the standard argument
\[
\min_{0\le k<T}\|\grad f(x_{k})\|_{*}
\le\sqrt{\frac{1}{T}\sum_{k=0}^{T-1}\|\grad f(x_{k})\|_{*}^{2}}
=O\!\big(T^{-1/2}\big).
\]

\section*{Special Cases}
\begin{itemize}
\item \textbf{Euclidean geometry.} $h(x)=\tfrac12\|x\|_{2}^{2}$ gives $\alpha=M=1$, $D_{h}(x,y)=\tfrac12\|x-y\|_{2}^{2}$. The bound reduces to the classic $O(1/\sqrt{T})$ rate for gradient descent on $L$‑smooth non‑convex $f$.
\item \textbf{Entropy prox.} For $h(x)=\sum_{i}x_{i}\log x_{i}$ on the simplex, $D_{h}$ becomes KL‑divergence. The theorem applies to problems such as logistic regression with probability simplex constraints, yielding the same $O(1/\sqrt{T})$ guarantee in the natural Bregman geometry.
\item \textbf{Adaptive stepsize.} If $\eta_{k}= \frac{1}{L+ \beta k}$ with $\beta>0$, the same telescoping argument yields a bound of order $O\big(\frac{\log T}{T}\big)$.
\end{itemize}

\end{document}